\section{Space Complexity --- Machines --- ARM}
%%% generated by the blueprint pipeline from TCSlib.Complexity.SpaceComplexity.Machines.ARM
\subsection*{Space Complexity --- Machines --- ARM: Introduction}

Declarations of the Lean module \texttt{TCSlib.Complexity.SpaceComplexity.Machines.ARM}.

\subsection*{Space Complexity --- Machines --- ARM: Results}

\begin{definition}[Logarithmic-program instructions]
\lean{Complexity.LogProg.Ins}
\label{Complexity.LogProg.Ins}
\leanok
\uses{Complexity.LogProg.Mode}
For natural numbers $m,d$ and a label type $\Lambda$, define instructions for a program
with $m$ registers and $d$ deciders, whose branch targets are labels in $\Lambda$. The
instructions increment a register, decrement it with truncated subtraction, clear it,
halve it, branch on whether it is zero or odd, branch on equality of two distinct
registers, call a decider on a virtual input specified by a mode and register arguments,
or return a Boolean answer and halt. They also include checks that the input has the
required unary and canonical binary encoding, branches testing whether a register's
binary expansion matches a specified input component, and branches on equality with the
first or second binary component of a nested input encoding.
\end{definition}

\begin{definition}[Abstract register machine]
\lean{Complexity.LogProg.ARM}
\label{Complexity.LogProg.ARM}
\leanok
\uses{Complexity.LogProg.Ins}
For natural numbers $m,d$ and a label type $\Lambda$, an abstract register machine is a
function assigning to each program label an instruction for a machine with $m$
registers, $d$ deciders, and labels in $\Lambda$.
\end{definition}

\begin{definition}[Plain input suffix]
\lean{Complexity.LogProg.plainWord}
\label{Complexity.LogProg.plainWord}
\leanok
\uses{Complexity.LogProg.tRun}
For a finite bit string $x$, let $t$ be the length of its initial run of $1$s. The plain
word of $x$ is the suffix obtained by deleting its first $t+2$ bits.
\end{definition}

\begin{definition}[Decoded pair of input words]
\lean{Complexity.LogProg.pairWords}
\label{Complexity.LogProg.pairWords}
\leanok
\uses{Complexity.LogProg.tRun, Turing.pairDecode}
For a finite bit string $x$, let $t$ be the length of its initial run of $1$s. Decode
the suffix of $x$ after its first $t+2$ bits using the aligned pair parser; if parsing
is undefined, assign the pair of empty words.
\end{definition}

\begin{lemma}[Recovery of the encoded payload]
\lean{Complexity.LogProg.plainWord_pairEncode}
\label{Complexity.LogProg.plainWord_pairEncode}
\leanok
\uses{Complexity.LogProg.dbl_replicate, Complexity.LogProg.plainWord, Complexity.LogProg.tRun_pairEncode, Turing.pairEncode, Turing.pairEncode_eq_dbl}
\difficulty{3}
For every natural number $n$ and bit string $w$, remove from the pairing of the unary
string of $n$ ones with $w$ the initial run of ones together with the following two
separator bits; the remaining string is exactly $w$.
\end{lemma}

\begin{lemma}[Decoding an encoded pair]
\lean{Complexity.LogProg.pairWords_pairEncode}
\label{Complexity.LogProg.pairWords_pairEncode}
\leanok
\uses{Complexity.LogProg.dbl_replicate, Complexity.LogProg.pairWords, Complexity.LogProg.tRun, Complexity.LogProg.tRun_pairEncode, Turing.pairDecode_pairEncode, Turing.pairEncode, Turing.pairEncode_eq_dbl}
\difficulty{4}
For every natural number $n$ and Boolean words $u,w$, remove the leading run of $1$s and
the next two bits from the encoding of the pair whose first component is $1^n$ and whose
second component is the encoding of $(u,w)$. Decoding the remaining word returns
$(u,w)$.
\end{lemma}

\begin{definition}[Abstract machine configuration]
\lean{Complexity.LogProg.AConf}
\label{Complexity.LogProg.AConf}
\leanok
For a natural number $m$ and a type of machine states $\Lambda$, an abstract
configuration consists of an optional state, a vector of $m$ natural-number register
values, and an optional Boolean result.
\end{definition}

\begin{definition}[Abstract machine step]
\lean{Complexity.LogProg.astep}
\label{Complexity.LogProg.astep}
\leanok
\uses{Complexity.LogProg.AConf, Complexity.LogProg.ARM, Complexity.LogProg.ValidPair, Complexity.LogProg.ValidPlain, Complexity.LogProg.callSegs, Complexity.LogProg.pairWords, Complexity.LogProg.plainWord, Complexity.LogProg.vword}
Let $m,d$ be natural numbers, let $\Lambda$ be a type of states, let $A$ assign an
instruction to each state, let an oracle assign a Boolean answer to each of $d$ deciders
and each finite bit string, and fix an input bit string $x$. One step leaves a halted
configuration unchanged; otherwise it executes the instruction at the current state:
increment, truncated-decrement, clear, or halve the selected register and move to the
indicated state; branch to one of two states according as a register is zero, odd, or
equal to another register; or call the selected oracle on the virtual input determined
by the mode, input, and register arguments, then branch according to its answer. A
return instruction halts with its specified Boolean result. A plain- or pair-validity
instruction continues to its specified state when the input has the corresponding valid
encoding, and otherwise halts with result false. The three input-equality instructions
branch according as the selected register's binary word equals, respectively, the plain
payload, the first decoded pair word, or the second decoded pair word.
\end{definition}

\begin{definition}[Abstract machine run]
\lean{Complexity.LogProg.arun}
\label{Complexity.LogProg.arun}
\leanok
\uses{Complexity.LogProg.AConf, Complexity.LogProg.ARM, Complexity.LogProg.astep}
Let $m,d$ be natural numbers, let $\Lambda$ be a type of states, let $A$ assign an
instruction to each state, let an oracle assign a Boolean answer to each of $d$ deciders
and each finite bit string, and fix an input bit string $x$. Starting from a
configuration $a$, the configuration after $n$ steps is the result of iterating the
one-step transition exactly $n$ times.
\end{definition}

\begin{definition}[Machine simulation phases]
\lean{Complexity.LogProg.Ph}
\label{Complexity.LogProg.Ph}
\leanok
Define the finite set of simulation phases to consist of the following distinct symbols:
start, incB, decL, decE, decB, clrE, h0, hF, hT, eqBy, eqBn, vU1, vS, vW0, vWF, vWT,
vrw1, vrw2, vP1($\ell$) for an optional Boolean $\ell$, vP2($b,\ell$) for a Boolean $b$
and optional Boolean $\ell$, jK2, jC, jRB($b$), jI1($b$), jI2($b$) for a Boolean $b$,
jP1, jP2F, jP2T, jD1, jD2F, and jD2T.
\end{definition}

\begin{definition}[Continue machine action]
\lean{Complexity.LogProg.goAct}
\label{Complexity.LogProg.goAct}
\leanok
\uses{Turing.Action}
For a natural number $m$ and a type of states $S$, define the action that continues in
state $s$ by leaving the input head and all $m$ work-tape heads stationary, making no
work-tape writes, emitting no output symbol, and setting the successor state to $s$.
\end{definition}

\begin{definition}[Return machine action]
\lean{Complexity.LogProg.retAct}
\label{Complexity.LogProg.retAct}
\leanok
\uses{Turing.Action}
For a natural number $m$, a type of states $S$, and a Boolean $b$, define the action
that halts and emits $b$ by leaving the input head and all $m$ work-tape heads
stationary, making no work-tape writes, emitting the Boolean symbol $b$, and having no
successor state.
\end{definition}

\begin{definition}[Second rightward scan action]
\lean{Complexity.LogProg.rw2Act}
\label{Complexity.LogProg.rw2Act}
\leanok
\uses{Complexity.LogProg.xAct, Turing.Action}
For a register index $r$ and states $s_0,n_x$, define an action depending on an optional
Boolean. When the Boolean is present, it is the action given by the leftward register
action with successor state $s_0$; when it is absent, it is the action given by the
rightward register action with successor state $n_x$.
\end{definition}

\begin{definition}[Default inactive action]
\lean{Complexity.LogProg.junkAct}
\label{Complexity.LogProg.junkAct}
\leanok
\uses{Turing.Action}
For $m$ work tapes over the Boolean alphabet and state set $S$, define the default
action to leave the input head and every work-tape head stationary, make no writes or
output, and halt.
\end{definition}

\begin{definition}[Instruction transition action]
\lean{Complexity.LogProg.insTr}
\label{Complexity.LogProg.insTr}
\leanok
\uses{Complexity.LogProg.Ins, Complexity.LogProg.Ph, Complexity.LogProg.backAct, Complexity.LogProg.clrEAct, Complexity.LogProg.cmpAct, Complexity.LogProg.dcmp1Act, Complexity.LogProg.dcmp2Act, Complexity.LogProg.decDAct, Complexity.LogProg.decEAct, Complexity.LogProg.decLAct, Complexity.LogProg.eqBAct, Complexity.LogProg.eqCAct, Complexity.LogProg.goAct, Complexity.LogProg.halfLAct, Complexity.LogProg.incBAct, Complexity.LogProg.incCAct, Complexity.LogProg.junkAct, Complexity.LogProg.pskip1Act, Complexity.LogProg.pskip2Act, Complexity.LogProg.retAct, Complexity.LogProg.rewAct, Complexity.LogProg.rw2Act, Complexity.LogProg.skipAct, Complexity.LogProg.toEndAct, Complexity.LogProg.valP1Act, Complexity.LogProg.valP2Act, Complexity.LogProg.valSAct, Complexity.LogProg.valUAct, Complexity.LogProg.valWAct, Complexity.LogProg.xAct, Turing.Action}
Given an instruction with $m$ registers and $d$ available deciders, a label, a phase, an
optional Boolean input symbol, and the current contents of each register, define the
corresponding one-step action of a machine whose states are pairs of labels and phases.
The action is determined by the instruction and phase, with the specified instruction
transitions and the default inactive action in the remaining cases.
\end{definition}

\begin{definition}[Labelled instruction transition action]
\lean{Complexity.LogProg.armTr}
\label{Complexity.LogProg.armTr}
\leanok
\uses{Complexity.LogProg.ARM, Complexity.LogProg.Ph, Complexity.LogProg.insTr, Turing.Action}
Let $A$ assign an instruction to each label in a set $\Lambda$. For a label, phase,
optional Boolean input symbol, and register contents, define the one-step action to be
the instruction transition action for the instruction assigned to that label, with
states given by pairs of labels and phases.
\end{definition}

\begin{definition}[Call specification for an instruction]
\lean{Complexity.LogProg.insCall}
\label{Complexity.LogProg.insCall}
\leanok
\uses{Complexity.LogProg.CallSpec, Complexity.LogProg.Ins, Complexity.LogProg.Ph}
Given an instruction with $m$ registers and $d$ available deciders and a phase, define
an optional call specification. It is present precisely when the instruction is a call
and the phase is the start phase; it records the chosen decider, input mode, argument
registers, and the positive and negative successor labels paired with the start phase.
In all other cases there is no call specification.
\end{definition}

\begin{definition}[Labelled call specification]
\lean{Complexity.LogProg.armCall}
\label{Complexity.LogProg.armCall}
\leanok
\uses{Complexity.LogProg.ARM, Complexity.LogProg.CallSpec, Complexity.LogProg.Ph, Complexity.LogProg.insCall}
Let $A$ assign an instruction to each label in $\Lambda$. For a label and phase, define
the optional call specification by taking the call specification of the assigned
instruction at that phase.
\end{definition}

\begin{definition}[Turing machine for a labelled program]
\lean{Complexity.LogProg.armProg}
\label{Complexity.LogProg.armProg}
\leanok
\uses{Complexity.LogProg.ARM, Complexity.LogProg.Ph, Complexity.LogProg.RProg, Complexity.LogProg.armCall, Complexity.LogProg.armTr}
Let $A$ assign an instruction to each label in $\Lambda$, and choose an initial label.
Define a machine with $m$ work tapes and states consisting of a label paired with a
phase. Its initial state is the chosen label paired with the start phase; its transition
function is the labelled instruction transition action, and its call function is the
labelled call specification.
\end{definition}

\begin{definition}[Initial configuration for a register computation]
\lean{Complexity.LogProg.aseam}
\label{Complexity.LogProg.aseam}
\leanok
\uses{Complexity.LogProg.Ph, Turing.Cfg, Turing.FinTM.bufferTape}
Given an input string $x$ of Boolean symbols, a label $l$, a natural-number value $v(r)$
for each of $m$ registers, and a Boolean output string, define a machine configuration
with state $(l,\mathrm{start})$, input head at the first position, each register tape
containing the binary expansion of its assigned value, every work-tape head at position
zero, and the specified output string.
\end{definition}
